\section{Secondary quantitative routes and local stability}
\label{sec:two-uniform}

This appendix derives discreteness and local quadratic stability from
two-uniformity alone, without the stabilizer hypothesis used by the main
cleaning argument.  It then shows how higher uniformity enlarges the
certified generator-coordinate region and strengthens the residual estimate
proved directly in Section~\ref{sec:quantitative}.  Readers interested only
in the headline exact and quantitative theorems may skip these comparisons.

\subsection{Discreteness from two-party maximal mixing}

The rigidity theorem uses the stabilizer hypothesis to identify the
symmetry group.  Finiteness of that group needs far less: two-party
marginals alone, at any local dimension, with no stabilizer structure.
Throughout this subsection \(q\geq2\) is an arbitrary integer; only the
statements naming Clifford operators return to \(q=p^e\).

Call a pure state
\(\lvert\psi\rangle\in(\mathbb C^q)^{\otimes n}\) \emph{\(2\)-uniform} if
\[
 \rho^\psi_{\{j,k\}}=q^{-2}I\qquad\text{for all }j\ne k .
\]
Tracing once more makes every one-party marginal \(q^{-1}I\).  Such states
exist only for \(n\geq4\), since at \(n=3\) purity forces
\(\rank\rho_{\{1,2\}}=\rank\rho_{\{3\}}\leq q\); and every
\(\AME(2m,q)\) state with \(m\geq2\) is \(2\)-uniform, stabilizer or not.
Write
\[
 G(\psi)=\{U_1\otimes\cdots\otimes U_n:
   (U_1\otimes\cdots\otimes U_n)\lvert\psi\rangle
     \in\mathbb C\lvert\psi\rangle\}
\]
for its product-unitary symmetry group and, for a product unitary \(U\),
\[
 \varepsilon(U)=\min_\theta\bigl\lVert U\lvert\psi\rangle
    -e^{i\theta}\lvert\psi\rangle\bigr\rVert
  =\bigl(2-2\lvert\langle\psi\rvert U\lvert\psi\rangle\rvert
   \bigr)^{1/2}
\]
for its defect.  Because operators on distinct sites commute, a product
unitary is a single exponential:
\(\bigotimes_je^{iH_j}=e^{iL}\) with \(L=\sum_jH_j^{(j)}\), where
\(H^{(j)}\) denotes \(H\) acting at site \(j\).  Splitting each \(H_j\)
into a traceless part \(h_j\) and a multiple of the identity leaves
\(L=M+cI\) with \(M=\sum_jh_j^{(j)}\).  Everything below rests on one
computation about \(M\).

\AMEFigureFour

\begin{lemma}[Local-generator isometry]
\label{lem:local-generator-isometry}
Let \(\lvert\psi\rangle\) be \(2\)-uniform and let \(h_j\) and \(h'_j\),
\(1\leq j\leq n\), be traceless Hermitian operators on \(\mathbb C^q\).
Put \(M=\sum_jh_j^{(j)}\) and \(M'=\sum_jh_j'^{(j)}\).  Then
\(\langle\psi\rvert M\lvert\psi\rangle=0\) and
\[
 \bigl\langle M\lvert\psi\rangle,\,M'\lvert\psi\rangle\bigr\rangle
  =\frac1q\sum_{j=1}^n\Tr(h_jh'_j).                      \tag{B.1}
\]
In particular \(\lVert M\lvert\psi\rangle\rVert^2
 =q^{-1}\sum_j\lVert h_j\rVert_F^2\), so
\((h_1,\ldots,h_n)\mapsto\sqrt q\,M\lvert\psi\rangle\) is an isometry from
the traceless local generators, carrying the product Frobenius form, into
the Hilbert space regarded as a real inner product space.
\end{lemma}

\begin{proof}
One-party maximal mixing gives \(\langle M\rangle=0\).  In the expansion of
\(\langle MM'\rangle\), equal-site terms are
\(\Tr(h_jh'_j)/q\), while distinct-site terms are
\(\Tr(h_j)\Tr(h'_k)/q^2=0\).  Thus the displayed inner product has the
stated real part.  Its imaginary part is \(\langle[M,M']\rangle/2i=0\),
because only traceless one-site commutators remain.
\end{proof}

Pairwise maximal mixing is essential: for a Bell pair,
\(M=h\otimes I-I\otimes h^T\) annihilates the state, producing the familiar
\(U\otimes\overline U\) continuous symmetry.

We also need the standard Lie-algebra form of a product-unitary generator.

\begin{lemma}[Product generators]\label{lem:product-lie}
Let \(\pi:U(q)^n\to U(q^n)\) send \((U_1,\ldots,U_n)\) to
\(U_1\otimes\cdots\otimes U_n\).  Then \(\pi\) is a smooth homomorphism of
compact Lie groups, its image \(P\) is closed, and
\[
 \operatorname{Lie}(P)
  =\Bigl\{i\textstyle\sum_{j}H_j^{(j)}:
    H_j=H_j^\dagger\ \text{on }\mathbb C^q\Bigr\}.       \tag{B.2}
\]
Thus every one-parameter subgroup of a closed subgroup of \(P\) has such a
generator.
\end{lemma}

\begin{proof}
The source is compact, so \(P\) is closed.  At the identity,
\[
 d\pi(iH_1,\ldots,iH_n)=i\sum_jH_j^{(j)}.
\]
Surjectivity onto \(P\) makes this map onto \(\operatorname{Lie}(P)\).
The final claim is the closed-subgroup theorem.
\end{proof}

\begin{theorem}[Discreteness from \(2\)-uniformity]
\label{thm:two-uniform-discrete}
Let \(\lvert\psi\rangle\) be a \(2\)-uniform pure state of \(n\) qudits of
local dimension \(q\).  Then \(\operatorname{Lie}(G(\psi))=i\mathbb RI\),
and \(G(\psi)/U(1)\cdot I\) is finite.  In particular every
\(\AME(2m,q)\) state with \(m\geq2\), stabilizer or not, has finite
product-unitary symmetry group modulo global phase.
\end{theorem}

\begin{proof}
The set \(\{U\in U(q^n):U\lvert\psi\rangle\in\mathbb C\lvert\psi\rangle\}\)
is a closed subgroup, so \(G(\psi)=P\cap\{\cdots\}\) is a closed subgroup
of \(P\) and hence a compact Lie group.  Let
\(X\in\operatorname{Lie}(G(\psi))\).  By Lemma~\ref{lem:product-lie},
\(X=iL\) with \(L=\sum_jH_j^{(j)}\) Hermitian; write \(L=M+cI\) as above.
Each \(e^{itL}\) fixes the ray of \(\lvert\psi\rangle\), so
\(e^{itL}\lvert\psi\rangle=\lambda(t)\lvert\psi\rangle\) with
\(\lambda(t)=\langle\psi\rvert e^{itL}\lvert\psi\rangle\) unimodular and
smooth in \(t\).  Differentiating at \(t=0\) shows that
\(\lvert\psi\rangle\) is an eigenvector of the Hermitian operator \(L\),
say \(L\lvert\psi\rangle=\mu\lvert\psi\rangle\) with \(\mu\) real.  Pairing
with \(\lvert\psi\rangle\) and using \(\langle M\rangle=0\) gives
\(\mu=c\), hence \(M\lvert\psi\rangle=0\).  By
Lemma~\ref{lem:local-generator-isometry}, \(\sum_j\lVert h_j\rVert_F^2=0\),
so every \(h_j=0\) and every \(H_j\) is scalar.  Thus
\(\operatorname{Lie}(G(\psi))=i\mathbb RI\).  Since \(U(1)\cdot I\) is a
closed central one-dimensional subgroup of \(G(\psi)\), the compact Lie
group \(G(\psi)/U(1)\cdot I\) has dimension zero and is therefore finite.
\end{proof}

\begin{remark}[Which torus is divided out]
\label{rem:local-phase-torus}
Upstairs in \(U(q)^n\) the preimage \(\pi^{-1}(G(\psi))\) always contains
the local-phase torus \(T=U(1)^n\), of dimension \(n\).  This does not
contradict the theorem: \(\pi(T)=U(1)\cdot I\) and \(\ker\pi\subset T\) has
dimension \(n-1\), so
\(\pi^{-1}(G(\psi))/T\cong G(\psi)/U(1)\cdot I\).  Statements about local
factors below are statements on \(\mathrm{PU}(q)^n\), not statements
modulo one global phase.
\end{remark}

\begin{remark}[Discreteness in the literature]
\label{rem:discreteness-prior-art}
Wirthm\"uller proves related finiteness criteria for binary stabilizer codes
\cite[Theorem~7 and Corollary~11]{Wirthmuller2011}; qubit stabilizer
symmetries are classified in \cite{EnglbrechtKraus2020}, and Tan computes
the four-qutrit AME group \cite[Theorem~5.3]{Tan2026}.  The theorem above
instead uses only \(2\)-uniformity and allows arbitrary local dimension and
nonstabilizer states.

The hypothesis cannot be reduced to one-party maximal mixing: such states
can have positive-dimensional local-unitary symmetry
\cite[Lemma~2]{SlowikSawickiMaciazek2021}.  Thus pairwise maximal mixing,
not merely criticality of the state, is what eliminates continuous product
symmetries.
\end{remark}

\begin{remark}[Two-uniformity does not identify the group]
\label{rem:two-uniform-not-clifford}
The equal-phase CSS state of \(C=\mathrm{RM}(1,4)\leq\F_2^{16}\) is
\(2\)-uniform because \(d(C)=8\) and \(d(C^\perp)=4\).  Yet
\(T^{\otimes16}\), \(T=\diag(1,e^{i\pi/4})\), fixes it: all codeword
weights are \(0,8,\) or \(16\).  Since \(T\) does not preserve Weyl axes,
two-uniformity forces finiteness but not Cliffordness.  The AME hypothesis
and Theorem~\ref{thm:lu-lc-rigidity} supply that second conclusion.
\end{remark}

\begin{corollary}[Discrete product-unitary symmetry]
\label{cor:discrete-lu-symmetry}
For a stabilizer AME state as in Theorem~\ref{thm:lu-lc-rigidity}, let
\(G\leq U(q)^{2m}\) be its local-factor tuple group and let
\(\widetilde G\leq U(q)^{2m}\rtimes S_{2m}\) also allow party
permutations.  The coordinatewise scalar torus \(T=(S^1)^{2m}\) is the
identity component of both groups, and
\(\Gamma=G/T\) and \(\widetilde\Gamma=\widetilde G/T\) are finite discrete
groups.  If \(\Pi\) is the realized party-permutation group, then
\[
  1\longrightarrow \Gamma\longrightarrow\widetilde\Gamma
    \longrightarrow\Pi\longrightarrow 1.
\]
This extension splits exactly when the realized permutations admit a
homomorphic choice of projective product-unitary lifts.
\end{corollary}

\begin{proof}
Theorem~\ref{thm:lu-lc-rigidity} puts every local factor in the finite
projective Clifford group.  Projectivization has kernel \(T\), so its image
is finite and discrete and \(T\) is the identity component.  Forgetting the
party permutation gives the third sequence.  Its splitting criterion is the
definition of a split extension.
\end{proof}


\subsection{Local quadratic stability}

Discreteness makes the exact symmetries isolated.  The same identity
\((B.1)\) says how fast the defect grows away from them, with a constant
that does not involve the number of parties.

\begin{theorem}[Local stability]\label{thm:two-uniform-stability}
Let \(\lvert\psi\rangle\) be \(2\)-uniform, let each \(h_j\) be traceless
Hermitian on \(\mathbb C^q\), and put
\[
 U=\bigotimes_{j=1}^ne^{ih_j},\qquad
 t=\sum_{j=1}^n\lVert h_j\rVert_{\mathrm{op}},\qquad
 D=\Bigl(\sum_{j=1}^n\lVert h_j\rVert_F^2\Bigr)^{1/2}.
\]
If \(t\leq\frac12\) then
\[
 \frac{D^2}q\Bigl(1-\frac t3\Bigr)\leq\varepsilon(U)^2
 \leq\frac{D^2}q\Bigl(1+\frac t3\Bigr).                  \tag{B.3}
\]
In particular \(D\leq\sqrt{6q/5}\,\varepsilon(U)<1.096\sqrt q\,
\varepsilon(U)\), and \(\varepsilon(U)^2/D^2\to q^{-1}\) as \(t\to0\),
uniformly in \(n\).
\end{theorem}

\begin{proof}
Since the sites commute, \(U=e^{iM}\) with
\(\lVert M\rVert_{\mathrm{op}}\leq t\).  Put
\(\mu_2=\langle M^2\rangle=D^2/q\), using
Lemma~\ref{lem:local-generator-isometry}.  Taylor's theorem and the spectral
calculus give
\[
 \langle e^{iM}\rangle=1-\mu_2/2+r,\qquad |r|\leq t\mu_2/6.
\]
Since \(\mu_2\leq t^2\leq1/4\), the real main term is positive, and
\[
 \Bigl\lvert\,\lvert\langle e^{iM}\rangle\rvert
   -(1-\mu_2/2)\Bigr\rvert\leq t\mu_2/6 .
\]
Substitution into
\(\varepsilon(U)^2=2-2\lvert\langle e^{iM}\rangle\rvert\) proves (B.3)
and its stated consequences.
\end{proof}

The constant \(\sqrt{6/5}=1.0954\ldots\) is therefore sharp to within about
ten percent: the true asymptotic constant is \(\sqrt q\), and the excess is
the cubic Taylor cutoff, not structure.

\begin{remark}[The stability constant is a Fisher metric]
\label{rem:fisher-isotropy}
For \(U(s)=e^{isM}\), the pure-state quantum Fisher information is
\(4\operatorname{Var}_\psi(M)\) \cite{BraunsteinCaves1994}.  Hence (B.1)
gives
\[
 F_Q=\frac4q\sum_{j=1}^n\lVert h_j\rVert_F^2 .
\]
Thus \(2\)-uniformity makes the Fisher form a flat multiple of the Euclidean
form on traceless one-site generators.  This determines the local constant,
but the size of the certified neighbourhood depends on higher moments and is
treated next.
\end{remark}

\subsection{Higher uniformity and the generator-coordinate region}

The size of the certified region is controlled by the state's uniformity
order, not by the number of parties.  Call
\(\lvert\psi\rangle\) \emph{\(k\)-uniform} if every reduced operator on at
most \(k\) parties is maximally mixed; \(2\)-uniformity is the case
\(k=2\), and an \(\AME(2m,q)\) state is \(m\)-uniform.  The hypothesis
\(t\leq\frac12\) in Theorem~\ref{thm:two-uniform-stability} pays for a
cubic Taylor remainder.  Each additional order of uniformity removes one
order from that remainder, because the state and the tracial state assign
the same value to every power of \(M\) below the uniformity order.

\begin{lemma}[Tracial agreement below the uniformity order]
\label{lem:tracial-agreement}
Let \(\lvert\psi\rangle\) be \(k\)-uniform on \(n\) qudits.  Then
\(\langle\psi\rvert O\lvert\psi\rangle=q^{-n}\Tr O\) for every operator
\(O\) supported on at most \(k\) parties.  In particular, with
\(M=\sum_jh_j^{(j)}\) as above,
\(\langle\psi\rvert M^r\lvert\psi\rangle=q^{-n}\Tr(M^r)\) for every
\(r\leq k\).
\end{lemma}

\begin{proof}
If \(O=O_S\otimes I\) with \(\lvert S\rvert\leq k\) then
\(\langle O\rangle=\Tr(\rho^\psi_SO_S)=q^{-\lvert S\rvert}\Tr O_S
=q^{-n}\Tr O\).  Expanding \(M^r\) gives a sum of products of one-site
operators touching at most \(r\leq k\) sites, so every term is of that
shape.
\end{proof}

The tracial state \(q^{-n}\Tr\) is the product of the one-site tracial
states, and the site summands of \(M\) commute, so
\(q^{-n}\Tr f(M)=\mathbb E\,f(X_1+\cdots+X_n)\) with the \(X_j\)
independent and \(X_j\) uniform on the spectrum of \(h_j\).  The tracial
main term is therefore a product of one-site characteristic functions, and
the following contraction estimate controls each factor.

\begin{lemma}[Per-site contraction]\label{lem:site-contraction}
Let \(h\) be traceless Hermitian on \(\mathbb C^q\) with spectral spread
\(s=\lambda_{\max}(h)-\lambda_{\min}(h)\).  Then
\[
 \bigl\lvert q^{-1}\Tr e^{ih}\bigr\rvert^2
   \leq1-c(s)\lVert h\rVert_F^2/q,
 \qquad
 c(s)=\max\Bigl\{1-\tfrac{s^2}{12},\ \tfrac4{\pi^2}\,[s\leq\pi]\Bigr\}.
\]
\end{lemma}

\begin{proof}
Write \(q^{-1}\Tr e^{ih}=\mathbb E\,e^{iX}\) with \(X\) uniform on the
spectrum.  For an independent copy \(X'\),
\(\lvert\mathbb E\,e^{iX}\rvert^2=\mathbb E\cos(X-X')\) and
\(\lvert X-X'\rvert\leq s\).  From
\(\cos u\leq1-u^2/2+u^4/24=1-(u^2/2)(1-u^2/12)\) and \(u^2\leq s^2\),
\(\mathbb E\cos(X-X')\leq1-(1-s^2/12)\operatorname{Var}X\); from
\(\cos u\leq1-2u^2/\pi^2\) for \(\lvert u\rvert\leq\pi\),
\(\mathbb E\cos(X-X')\leq1-(4/\pi^2)\operatorname{Var}X\).  Tracelessness
gives \(\operatorname{Var}X=\Tr(h^2)/q=\lVert h\rVert_F^2/q\).
\end{proof}

\begin{theorem}[Stability at uniformity order \(k\)]
\label{thm:k-uniform-stability}
Let \(\lvert\psi\rangle\) be \(k\)-uniform with \(k\geq2\), let each
\(h_j\) be traceless Hermitian with spectral spread at most \(s\), and keep
\(t\), \(D\), \(U\) as in Theorem~\ref{thm:two-uniform-stability}.  Put
\[
 c=c(s),\qquad
 \lambda=\frac{2t^{\,k-1}}{(k+1)!},\qquad
 \rho=c-\frac{c^2D^2}{4q}-2\lambda .
\]
Then \(\varepsilon(U)^2\geq\rho D^2/q\); in particular \(\rho>0\) gives
\(D\leq\sqrt{q/\rho}\;\varepsilon(U)\).
\end{theorem}

\begin{proof}
Write \(e^{ix}=\sum_{r\leq k}(ix)^r/r!+\varrho(x)\) with
\(\lvert\varrho(x)\rvert\leq\lvert x\rvert^{k+1}/(k+1)!\).  By
Lemma~\ref{lem:tracial-agreement} the state and the tracial state agree on
every \(M^r\) with \(r\leq k\), so
\(\langle e^{iM}\rangle-q^{-n}\Tr e^{iM}
=\langle\varrho(M)\rangle-q^{-n}\Tr\varrho(M)\).  For any spectral weight
vector \(p_i\geq0\) with \(\sum_ip_i=1\),
\[
 \Bigl\lvert\sum_ip_i\varrho(\lambda_i)\Bigr\rvert
 \leq\frac{\lVert M\rVert_{\mathrm{op}}^{\,k-1}}{(k+1)!}\sum_ip_i\lambda_i^2 ,
\]
and both states give \(M^2\) the value \(D^2/q\), by
Lemma~\ref{lem:tracial-agreement} with \(r=2\) on one side and by
tracelessness on the other.  Since
\(\lVert M\rVert_{\mathrm{op}}\leq t\), each term is at most
\(t^{k-1}D^2/(q(k+1)!)\), so the two expectations differ by at most
\(\lambda D^2/q\).  The sites commute, so
\(q^{-n}\Tr e^{iM}=\prod_jq^{-1}\Tr e^{ih_j}\), and
Lemma~\ref{lem:site-contraction} with \(1-y\leq e^{-y}\) gives
\[
 \bigl\lvert\langle e^{iM}\rangle\bigr\rvert
  \leq\prod_j\bigl(1-c\lVert h_j\rVert_F^2/q\bigr)^{1/2}
      +\lambda D^2/q
  \leq e^{-cD^2/(2q)}+\lambda D^2/q .
\]
With \(x=cD^2/(2q)\) and \(1-e^{-x}\geq x-x^2/2\),
\(\varepsilon(U)^2=2-2\lvert\langle e^{iM}\rangle\rvert
\geq2x-x^2-2\lambda D^2/q=\rho D^2/q\).
\end{proof}

\begin{corollary}[The certified radius grows with the uniformity order]
\label{cor:k-uniform-region}
Let \(\lvert\psi\rangle\) be \(k\)-uniform with \(k\geq2\).  If every
\(h_j\) has spectral spread at most \(\frac12\), if \(D\leq\sqrt q/2\), and
if
\[
 t\leq R_k:=\Bigl(\frac{(k+1)!}{48}\Bigr)^{1/(k-1)},
\]
then \(D\leq\sqrt{6q/5}\,\varepsilon(U)\).  One has
\(\log R_k=\log k-1+O(\log k/k)\), so \(R_k\sim k/e\).
\end{corollary}

\begin{proof}
Spread at most \(\frac12\) gives \(c\geq1-\frac1{48}=0.97917\); with
\(D^2\leq q/4\), \(c^2D^2/(4q)\leq c^2/16\), so
\(c-c^2D^2/(4q)\geq0.91924\); and \(t^{k-1}\leq(k+1)!/48\) gives
\(2\lambda\leq\frac1{12}\).  Hence \(\rho\geq0.83591>\frac56\) and
\(\sqrt{q/\rho}\leq\sqrt{6q/5}\).  The asymptotics are Stirling's formula
applied to \(\log R_k=(\log(k+1)!-\log48)/(k-1)\).
\end{proof}

At \(k=2\) the direct theorem is stronger; for an \(\AME(2m,q)\) state,
\(k=m\) and \(R_k\sim k/e\), so the certified generator radius grows
linearly in the party count.  The following scalar example shows that no
uniform estimate of this form can grow faster.

\begin{proposition}[Ceiling on the certified radius]
\label{prop:region-ceiling}
For every \(q\geq2\) and every \(n\) there are traceless Hermitian \(h_j\)
on \(\mathbb C^q\) with
\(\sum_j\lVert h_j\rVert_{\mathrm{op}}=2\pi(1-1/q)n\) and \(D>0\) for which
\(\bigotimes_je^{ih_j}\) is a scalar, hence \(\varepsilon=0\).  No estimate
\(\varepsilon(U)^2\geq\rho D^2/q\) with \(\rho>0\) can therefore hold on a
region of radius exceeding \(2\pi(1-1/q)n\).
\end{proposition}

\begin{proof}
Take \(h_j=(2\pi/q)\diag(q-1,-1,\ldots,-1)\) at every site.  It is
traceless, and both eigenvalues of \(e^{ih_j}\) equal \(e^{-2\pi i/q}\), so
\(e^{ih_j}=e^{-2\pi i/q}I\) and the product of \(n\) of them is a global
phase.  Here \(\lVert h_j\rVert_{\mathrm{op}}=2\pi(q-1)/q\) and
\(\lVert h_j\rVert_F^2=(2\pi/q)^2q(q-1)>0\).
\end{proof}

Thus linear growth is the maximal possible order.  At fixed uniformity
order the next family gives a length-independent obstruction.

\begin{proposition}[The certified radius at uniformity order three]
\label{prop:stability-region}
Let \(\ell\geq3\), \(N=2^\ell\), and let \(\lvert C\rangle\) be the
equal-phase CSS state of \(C=\mathrm{RM}(1,\ell)\leq\F_2^N\).  Its
uniformity order is \(3\), for every \(\ell\).  Take
\(h_j=cZ(1)\) at every site and put \(\theta=Nc\in(0,2\pi)\).  Then
\[
 \sum_j\lVert h_j\rVert_{\mathrm{op}}=\theta,\qquad
 D=\frac{\theta\sqrt2}{\sqrt N},\qquad
 \varepsilon=\frac{2\lvert\sin(\theta/2)\rvert}{\sqrt N},\qquad
 \frac D{\sqrt2\,\varepsilon}=\frac\theta{2\sin(\theta/2)} .
\]
The last ratio does not depend on \(N\) and increases on \((0,2\pi)\); it
passes \(\sqrt{6/5}\) at the unique \(\theta^*=1.46562\ldots\).  Hence the
conclusion of Theorem~\ref{thm:two-uniform-stability} fails for every
\(t>\theta^*\), for every length.  At uniformity order three the supremal
radius for the constant \(\sqrt{6q/5}\) therefore lies between
\(R_3=1/\sqrt2\) and \(\theta^*\).
\end{proposition}

\begin{proof}
The nonzero weights of \(\mathrm{RM}(1,\ell)\) are \(N/2\) and \(N\), and
\(C^\perp=\mathrm{RM}(\ell-2,\ell)\) has minimum distance \(4\).  A
marginal of the equal-phase CSS state fails to be maximally mixed only when
a nonzero word of \(C\) or of \(C^\perp\) is supported inside it, so by
\((2.3)\) the uniformity order is
\(\min(N/2,4)-1=3\) whenever \(N\geq8\).  Since
\(Z(1)\lvert x\rangle=(-1)^x\lvert x\rangle\) at a single site,
\(U\lvert x\rangle=e^{ic(N-2\operatorname{wt}(x))}\lvert x\rangle\), and
the weight distribution \(1,\,2N-2,\,1\) on weights \(0,\,N/2,\,N\) gives
\(\langle U\rangle=1-(1-\cos\theta)/N\).  This is nonnegative, so
\(\varepsilon^2=2-2\langle U\rangle=4\sin^2(\theta/2)/N\).  Also
\(\lVert cZ(1)\rVert_F^2=2c^2\) and
\(\lVert cZ(1)\rVert_{\mathrm{op}}=c\), giving the displayed \(D\) and
\(\theta\).  Finally \(\theta/(2\sin(\theta/2))=u/\sin u\) with
\(u=\theta/2\in(0,\pi)\), which increases from \(1\) to \(\infty\); for
every \(\theta>\theta^*\) the ratio exceeds the constant of
Theorem~\ref{thm:two-uniform-stability}.  The violation begins at the
limiting defect \(2\sin(\theta^*/2)N^{-1/2}\).
\end{proof}

\AMEFigureFive

The example fixes the relevant parameter: the obstruction remains at
uniformity order three as \(N\) grows, whereas an \(\AME(2m,q)\) state has
order \(m\).  These are generator-coordinate estimates for known product
unitaries, not certification statements for an unknown device.

\paragraph{Comparison with the balanced-cut estimate.}
Proposition~\ref{prop:main-residual-stability} removes the global
\(\ell^1\) budget, permits per-site spectral spread up to \(\pi\), and gives
\[
 D\leq2\sqrt{q/c(s)}\,\varepsilon(U),
 \qquad
 c(s)=\max\left\{1-\frac{s^2}{12},\frac4{\pi^2}\right\}.
\]
It is complementary to the moment estimate above and is proved in the main
text.

\subsection{From local stability to exact branch selection}

\begin{corollary}[Approximate symmetries decompose]
\label{cor:approximate-decomposition}
Let \(\lvert\psi\rangle\) be a stabilizer \(\AME(2m,q)\) state with
\(m\geq2\).  There is an \(\varepsilon_0(\psi)>0\) such that every product
unitary \(U\) with \(\varepsilon(U)<\varepsilon_0\) factors as
\[
 U=g\cdot\bigotimes_{j=1}^{2m}e^{ih_j},\qquad g\in G(\psi),
\]
with the \(h_j\) traceless Hermitian,
\(\sum_j\lVert h_j\rVert_{\mathrm{op}}\leq\frac12\), and
\[
 \Bigl(\sum_j\lVert h_j\rVert_F^2\Bigr)^{1/2}
   \leq\sqrt{6q/5}\;\varepsilon(U).
\]
By Theorem~\ref{thm:lu-lc-rigidity} the factor \(g\) is an exact
local-Clifford symmetry.
\end{corollary}

The size constraint on the \(h_j\) is part of the conclusion and cannot be
dropped: the generators of Proposition~\ref{prop:region-ceiling} are
traceless with \(\bigotimes_je^{ih_j}\) scalar, so the defect vanishes
while the generator norm does not.

\begin{proof}
On \(X=\mathrm{PU}(q)^{2m}\), put \(f(U)=\varepsilon(U)^2\).  Let
\(B\subseteq X\) be the image of
\(\{(h_j)\ \text{traceless}:\sum_j\lVert h_j\rVert_{\mathrm{op}}
\leq\frac12\}\) under \((h_j)\mapsto[\bigotimes_je^{ih_j}]\).  That map is
a local diffeomorphism at the origin, so \(B\) is a neighbourhood of the
identity coset.  Theorem~\ref{thm:two-uniform-stability} gives
\[
 f\geq\frac5{6q}\sum_j\lVert h_j\rVert_F^2
\]
there.  Since \(f(gV)=f(V)\) for every exact symmetry \(g\), the same
estimate holds on every translate \(gB\).

The zero set \(Z\) of \(f\) is finite by
Theorem~\ref{thm:two-uniform-discrete} and
Remark~\ref{rem:local-phase-torus}.  So \(\bigcup_{g\in Z}gB\) is a
neighbourhood of \(Z\).  On its compact complement \(K\), either
\(K=\varnothing\) or \(\min_Kf>0\).  Taking
\(\varepsilon_0^2<\min_Kf\) in the latter case forces \(U\in gB\) for some
\(g\in Z\), and the translated quadratic estimate gives the conclusion.
\end{proof}

The factorization holds for every \(2\)-uniform state; the stabilizer
hypothesis identifies \(g\) as Clifford.  Its compactness threshold is not
explicit.  Section~\ref{sec:quantitative} supplies an explicit cleaning
radius, and Corollary~\ref{cor:relative-intertwiner-rounding} transfers it
without loss when an exact base intertwiner is known.

The same conclusion has the standard gate-theoretic interpretation.  If a
two-qudit gate \(W\) reshapes to an
\(\AME(4,q)\) state, then reshaping identifies its local equivalence group
with the product-unitary symmetry group of that state: an action on the two
input legs becomes complex conjugation on the corresponding state factors.
Theorem~\ref{thm:two-uniform-discrete} and
Remark~\ref{rem:local-phase-torus} therefore make the local equivalence group
finite modulo one-site phases.  For gates arising from stabilizer
\(\AME(4,q)\) states it is locally Clifford, by
Theorem~\ref{thm:lu-lc-rigidity}.  This is the usual two-unitary
correspondence; see
\cite{RatherAravindaLakshminarayan2022,Pahari2026}.
